\documentclass[12pt,a4paper]{report}

% ── Encodage et langue ────────────────────────────────────────────────────────
\usepackage[utf8]{inputenc}
\usepackage[T1]{fontenc}
\usepackage[french]{babel}

% ── Mise en page ──────────────────────────────────────────────────────────────
\usepackage[top=2.5cm, bottom=2.5cm, left=3cm, right=2.5cm]{geometry}
\usepackage{setspace}
\onehalfspacing

% ── Typographie ───────────────────────────────────────────────────────────────
\usepackage{lmodern}
\usepackage{microtype}

% ── Graphiques et couleurs ────────────────────────────────────────────────────
\usepackage{xcolor}
\usepackage{graphicx}
\usepackage{tikz}
\usetikzlibrary{shapes.geometric, arrows.meta, positioning, fit, backgrounds}

% ── Tableaux ──────────────────────────────────────────────────────────────────
\usepackage{booktabs}
\usepackage{tabularx}
\usepackage{multirow}
\usepackage{array}

% ── Code source ───────────────────────────────────────────────────────────────
\usepackage{listings}
\usepackage{inconsolata}

% ── Maths ─────────────────────────────────────────────────────────────────────
\usepackage{amsmath}
\usepackage{amssymb}

% ── Liens et PDF ──────────────────────────────────────────────────────────────
\usepackage[hidelinks]{hyperref}

% ── En-têtes et pieds de page ─────────────────────────────────────────────────
\usepackage{fancyhdr}
\pagestyle{fancy}
\fancyhf{}
\fancyhead[L]{\small RecoShop}
\fancyhead[R]{\small \leftmark}
\fancyfoot[C]{\thepage}
\renewcommand{\headrulewidth}{0.4pt}

% ── Couleurs projet ───────────────────────────────────────────────────────────
\definecolor{gold}{RGB}{212,175,55}
\definecolor{darkbg}{RGB}{26,26,46}
\definecolor{codebg}{RGB}{245,245,245}
\definecolor{codecomment}{RGB}{100,100,100}
\definecolor{codestring}{RGB}{163,21,21}
\definecolor{codekeyword}{RGB}{0,0,200}

% ── Style de code ─────────────────────────────────────────────────────────────
\lstset{
  backgroundcolor=\color{codebg},
  basicstyle=\footnotesize\ttfamily,
  keywordstyle=\color{codekeyword}\bfseries,
  stringstyle=\color{codestring},
  commentstyle=\color{codecomment}\itshape,
  breaklines=true,
  showstringspaces=false,
  frame=single,
  rulecolor=\color{gray!40},
  numbers=left,
  numberstyle=\tiny\color{gray},
  numbersep=8pt,
  tabsize=4,
  captionpos=b,
}

\lstdefinelanguage{Python}{
  keywords={def,class,return,if,else,elif,for,while,import,from,as,try,except,
             in,not,and,or,True,False,None,with,lambda,pass,break,continue,raise},
  morestring=[b]",
  morestring=[b]',
  morecomment=[l]{\#},
  sensitive=true,
}

% ── Commandes personnalisées ───────────────────────────────────────────────────
\newcommand{\fichier}[1]{\texttt{#1}}
\newcommand{\fonction}[1]{\texttt{#1()}}
\newcommand{\classe}[1]{\texttt{#1}}
\newcommand{\var}[1]{\texttt{#1}}

% ── Titre ─────────────────────────────────────────────────────────────────────
\title{
  \vspace{2cm}
  {\color{gold}\rule{\linewidth}{1.5pt}}\\[0.4cm]
  {\Huge\bfseries RecoShop}\\[0.3cm]
  {\large Systeme de recommandation e-commerce}\\[0.2cm]
  {\large Application web Django avec moteur hybride}\\[0.4cm]
  {\color{gold}\rule{\linewidth}{1.5pt}}\\[2cm]
}
\author{
  Rapport de projet\\[0.5cm]
  \large Universite\\[0.2cm]
}
\date{\today}

% =============================================================================
\begin{document}
% =============================================================================

\maketitle
\thispagestyle{empty}

\newpage
\tableofcontents
\newpage

% =============================================================================
\chapter{Introduction}
% =============================================================================

\section{Presentation generale}

RecoShop est une application web e-commerce developpee avec le framework Django.
Son objectif principal est de proposer a chaque utilisateur des recommandations
de produits personnalisees, en combinant plusieurs algorithmes d'apprentissage
automatique au sein d'un moteur hybride.

L'application couvre la totalite du cycle de vie d'une boutique en ligne : 
inscription et authentification des utilisateurs, catalogue de produits avec 
recherche et filtrage, panier d'achat, finalisation des commandes, gestion du 
profil, et page de recommandations dediee. Le systeme de recommandation 
constitue la valeur ajoutee principale du projet d'un point de vue technique 
et academique.

\section{Objectifs du projet}

Les objectifs poursuivis sont les suivants :

\begin{itemize}
  \item Concevoir une boutique en ligne fonctionnelle et deployable en production.
  \item Implementer trois familles d'algorithmes de recommandation : filtrage 
        collaboratif base utilisateur, filtrage collaboratif base article, et 
        filtrage par contenu.
  \item Combiner ces algorithmes dans un systeme hybride avec gestion automatique 
        du cold start.
  \item Evaluer la qualite des recommandations avec des metriques standards : 
        Precision@K, Recall@K, F1-Score et NDCG@K.
  \item Deployer l'application sur une plateforme cloud (Render) avec base de 
        donnees PostgreSQL.
\end{itemize}

\section{Technologies utilisees}

\begin{table}[h]
\centering
\begin{tabularx}{\linewidth}{@{}lX@{}}
\toprule
\textbf{Composant} & \textbf{Technologie} \\
\midrule
Framework web       & Django 5.2 (Python) \\
Base de donnees     & SQLite (developpement) / PostgreSQL (production) \\
ORM                 & Django ORM avec \texttt{dj-database-url} \\
Fichiers statiques  & WhiteNoise avec compression \\
Machine Learning    & scikit-learn 1.8, NumPy 2.4, SciPy 1.17 \\
Serveur production  & Gunicorn + Render Web Service \\
Email               & SMTP Gmail via Django \texttt{send\_mail} \\
Images              & Pillow 12, URLs Unsplash \\
Fuseaux horaires    & \texttt{Africa/Dakar}, \texttt{tzdata} \\
\bottomrule
\end{tabularx}
\caption{Stack technologique de RecoShop}
\end{table}

% =============================================================================
\chapter{Architecture du projet}
% =============================================================================

\section{Structure des fichiers}

Le projet suit la structure conventionnelle d'une application Django
mono-application. Le dossier racine contient le package de configuration
\fichier{ecommerce\_reco/} et l'application principale \fichier{recommendations/}.

\begin{lstlisting}[language={},caption={Arborescence du projet},label={lst:tree}]
ecommerce_reco/
|-- manage.py
|-- requirements.txt
|-- build.sh
|-- .env
|-- ecommerce_reco/
|   |-- settings.py
|   |-- urls.py
|   |-- asgi.py
|   `-- wsgi.py
`-- recommendations/
    |-- models.py
    |-- views.py
    |-- urls.py
    |-- collaborative.py
    |-- content_based.py
    |-- hybrid.py
    |-- evaluation.py
    |-- admin.py
    |-- apps.py
    |-- context_processors.py
    |-- management/
    |   `-- commands/
    |       |-- seed_data.py
    |       `-- fetch_images.py
    |-- migrations/
    |   |-- 0001_initial.py
    |   `-- 0002_add_category_image_url.py
    |-- static/
    |   `-- recommendations/
    |       |-- css/main.css
    |       `-- js/main.js
    `-- templates/
        `-- recommendations/
            |-- base.html
            |-- home.html
            |-- products.html
            |-- product_detail.html
            |-- cart.html
            |-- profile.html
            |-- recommendations.html
            |-- login.html
            |-- register.html
            |-- landing.html
            |-- email_sent.html
            |-- activation_invalid.html
            |-- emails/
            |   `-- activation.html
            `-- partials/
                `-- product_card.html
\end{lstlisting}

\section{Roles des fichiers principaux}

\begin{table}[h]
\centering
\begin{tabularx}{\linewidth}{@{}lX@{}}
\toprule
\textbf{Fichier} & \textbf{Role} \\
\midrule
\fichier{models.py}         & Definition des modeles de donnees (7 classes) \\
\fichier{views.py}          & Logique de traitement des requetes (18 vues) \\
\fichier{urls.py}           & Routage des URLs vers les vues \\
\fichier{collaborative.py}  & Algorithmes de filtrage collaboratif \\
\fichier{content\_based.py} & Algorithme de filtrage par contenu TF-IDF \\
\fichier{hybrid.py}         & Orchestration du systeme hybride et cold start \\
\fichier{evaluation.py}     & Calcul des metriques d'evaluation \\
\fichier{settings.py}       & Configuration globale Django \\
\fichier{context\_processors.py} & Injection du compteur panier dans tous les templates \\
\fichier{seed\_data.py}     & Commande de population initiale de la base \\
\bottomrule
\end{tabularx}
\caption{Roles des fichiers principaux}
\end{table}

\section{Configuration et deploiement}

\subsection{Variables d'environnement}

Toute la configuration sensible est externalisee dans un fichier \fichier{.env}
charge par \texttt{python-dotenv}. Les variables cles sont les suivantes :
\var{SECRET\_KEY}, \var{DEBUG}, \var{ALLOWED\_HOSTS}, \var{DATABASE\_URL},
\var{CSRF\_TRUSTED\_ORIGINS}, \var{EMAIL\_HOST\_USER} et
\var{EMAIL\_HOST\_PASSWORD}.

\subsection{Base de donnees}

La configuration de la base de donnees utilise \texttt{dj-database-url} qui
parse automatiquement la variable \var{DATABASE\_URL}. En l'absence de cette
variable, elle bascule sur SQLite local :

\begin{lstlisting}[language=Python,caption={Configuration de la base de donnees}]
DATABASES = {
    'default': dj_database_url.config(
        default=f"sqlite:///{BASE_DIR / 'db.sqlite3'}",
        conn_max_age=600,
    )
}
\end{lstlisting}

\subsection{Script de build pour Render}

Le fichier \fichier{build.sh} est execute automatiquement a chaque deploiement
sur la plateforme Render :

\begin{lstlisting}[language=bash,caption={Script de build}]
pip install -r requirements.txt
python manage.py collectstatic --no-input
python manage.py migrate
\end{lstlisting}

Les fichiers statiques sont servis par WhiteNoise avec compression et hachage
de contenu en production (\texttt{CompressedManifestStaticFilesStorage}).

% =============================================================================
\chapter{Modeles de donnees}
% =============================================================================

\section{Vue d'ensemble}

L'application repose sur sept modeles Django. La figure suivante illustre leurs
relations.

\begin{figure}[h]
\centering
\begin{tikzpicture}[
  entity/.style={rectangle, draw=gold, fill=darkbg!10, thick,
                 rounded corners=3pt, minimum width=2.8cm, minimum height=0.9cm,
                 font=\small\bfseries},
  rel/.style={-Stealth, thick, color=gray!60},
  note/.style={font=\tiny, color=gray}
]
  \node[entity] (user)     at (0,0)     {CustomUser};
  \node[entity] (product)  at (6,0)     {Product};
  \node[entity] (category) at (6,3)     {Category};
  \node[entity] (rating)   at (3,-2.5)  {Rating};
  \node[entity] (comment)  at (3,-4.5)  {Comment};
  \node[entity] (cart)     at (0,-2.5)  {CartItem};
  \node[entity] (purchase) at (6,-2.5)  {Purchase};

  \draw[rel] (category) -- node[note,right]{1..N} (product);
  \draw[rel] (user)     -- node[note,above]{1..N} (rating);
  \draw[rel] (product)  -- node[note,right]{1..N} (rating);
  \draw[rel] (user)     -- node[note,left]{1..N}  (cart);
  \draw[rel] (product)  -- node[note,below]{1..N} (cart);
  \draw[rel] (user)     -- node[note,above]{1..N} (comment);
  \draw[rel] (product)  -- node[note,right]{1..N} (comment);
  \draw[rel] (user)     -- node[note,above]{1..N} (purchase);
  \draw[rel] (product)  -- node[note,right]{1..N} (purchase);
\end{tikzpicture}
\caption{Diagramme relationnel simplifie des modeles}
\end{figure}

\section{Description des modeles}

\subsection{CustomUser}

Etend \classe{AbstractUser} de Django. Le champ d'authentification est
remplace par l'adresse email (au lieu du nom d'utilisateur par defaut).
Les champs supplementaires sont \var{profile\_pic} (ImageField),
\var{bio} (TextField limite a 300 caracteres) et \var{date\_joined}
(auto-rempli a la creation).

\begin{lstlisting}[language=Python,caption={Modele CustomUser (extrait)}]
class CustomUser(AbstractUser):
    email = models.EmailField(unique=True)
    profile_pic = models.ImageField(upload_to='profile_pics/', null=True, blank=True)
    bio = models.TextField(max_length=300, blank=True)
    USERNAME_FIELD = 'email'
    REQUIRED_FIELDS = ['username']
\end{lstlisting}

\subsection{Category}

Represente une categorie de produits avec un slug unique utilise dans les URLs,
une icone (caractere emoji), une couleur de fond et une URL d'image optionnelle.
Le projet seed 12 categories : Electronique, Mode, Maison, Sport, Beaute,
Livres, Cuisine, Jeux, Voyage, Auto, Bijoux et Bebe.

\subsection{Product}

Modele central de l'application. Les champs notables sont :

\begin{itemize}
  \item \var{price} et \var{original\_price} : permettent de calculer le
        pourcentage de reduction via la methode \fonction{discount\_percent}.
  \item \var{tags} : chaine de caracteres separee par des virgules,
        utilisee pour la recherche et le filtrage TF-IDF.
  \item \var{is\_featured} : booleen pour mettre en avant certains produits.
  \item \var{image\_url} : URL Unsplash pour l'affichage des images.
\end{itemize}

Trois methodes de calcul sont definies directement sur le modele :
\fonction{average\_rating}, \fonction{rating\_count}, et
\fonction{discount\_percent}.

\subsection{Rating}

Lie un \classe{CustomUser} a un \classe{Product} avec un score entier entre 1
et 5. La contrainte \var{unique\_together} garantit qu'un utilisateur ne peut
noter un produit qu'une seule fois. C'est la donnee fondamentale exploitee par
le moteur de recommandation collaboratif.

\subsection{Comment}

Commentaire texte libre (max 1000 caracteres) d'un utilisateur sur un produit,
independant de la note.

\subsection{CartItem}

Represente une ligne du panier actif d'un utilisateur. La methode
\fonction{subtotal} calcule le prix total de la ligne
(\var{price} x \var{quantity}). La contrainte \var{unique\_together} sur
(\var{user}, \var{product}) garantit l'unicite de chaque produit dans le panier.

\subsection{Purchase}

Enregistre les commandes passees. Lors du checkout, chaque \classe{CartItem}
est converti en \classe{Purchase}, puis le panier est vide.

% =============================================================================
\chapter{Fonctionnalites de l'application}
% =============================================================================

\section{Authentification}

\subsection{Inscription avec confirmation par email}

Le processus d'inscription est le suivant : le formulaire collecte l'email,
le nom d'utilisateur et le mot de passe (avec confirmation). Si toutes les
validations passent, le compte est cree avec \var{is\_active = False}. Un token
signe est genere via \texttt{default\_token\_generator} de Django, encode en
base64, et envoye par email SMTP Gmail dans un template HTML
\fichier{emails/activation.html}.

Lors du clic sur le lien d'activation, la vue \fonction{activate\_view}
verifie la validite du token, active le compte et connecte l'utilisateur.
En cas d'echec SMTP, le compte est active directement et l'utilisateur est
connecte avec un message d'avertissement (mode de secours).

\begin{lstlisting}[language=Python,caption={Generation du lien d'activation}]
uid   = urlsafe_base64_encode(force_bytes(user.pk))
token = default_token_generator.make_token(user)
activation_link = f"{protocol}://{domain}/activate/{uid}/{token}/"
\end{lstlisting}

\subsection{Connexion et deconnexion}

La connexion utilise \texttt{authenticate()} avec l'email comme identifiant.
La redirection post-connexion respecte le parametre \var{next} de l'URL si
present. La deconnexion est protegee par \texttt{@login\_required} pour eviter
les requetes GET malveillantes.

\section{Catalogue de produits}

La vue \fonction{products\_view} offre les fonctionnalites suivantes.

La \textbf{recherche full-text} filtre simultanement sur le nom, la description,
les tags et la marque via des requetes \texttt{Q} Django combinant quatre
\texttt{icontains}.

Le \textbf{filtrage par categorie} est effectue via le slug de la categorie
passe en parametre GET (\var{?category=electronique}).

Le \textbf{tri} propose cinq options : produits vedettes, meilleures notes,
prix croissant, prix decroissant, et plus recents.

Les annotations Django (\texttt{Avg}, \texttt{Count}) calculent la note moyenne
et le nombre de notes directement en base de donnees, sans N+1 requetes.

\section{Detail d'un produit}

La page detail gere plusieurs actions dans une seule vue :

\begin{itemize}
  \item Affichage du produit avec ses commentaires et sa note moyenne.
  \item Soumission d'une note (1 a 5) via \texttt{update\_or\_create}.
  \item Ajout d'un commentaire texte.
  \item Affichage de 6 produits similaires calcules par filtrage par contenu.
  \item Detection de la presence du produit dans le panier de l'utilisateur.
\end{itemize}

\section{Panier et commandes}

Le panier supporte les requetes AJAX pour une mise a jour sans rechargement de
page. Les vues \fonction{add\_to\_cart} et \fonction{remove\_from\_cart}
detectent les requetes asynchrones via l'en-tete \texttt{X-Requested-With} et
retournent une reponse JSON le cas echeant.

Le checkout transforme chaque \classe{CartItem} en \classe{Purchase} dans une
boucle transactionnelle, puis supprime tous les items du panier.

\begin{lstlisting}[language=Python,caption={Logique du checkout}]
for item in cart_items:
    Purchase.objects.create(
        user=request.user,
        product=item.product,
        quantity=item.quantity
    )
cart_items.delete()
\end{lstlisting}

\section{Profil utilisateur}

La page profil affiche l'historique des achats, l'historique des notes, et
8 recommandations personnalisees. Si l'utilisateur a note au moins 3 produits,
les metriques d'evaluation du moteur sont affichees. La mise a jour du profil
permet de changer le nom d'utilisateur, la biographie, la photo de profil, et
le mot de passe. En cas de changement de mot de passe, \texttt{update\_session\_auth\_hash}
est appele pour eviter la deconnexion automatique.

\section{Context processor}

Le fichier \fichier{context\_processors.py} injecte automatiquement dans
chaque template le nombre d'articles dans le panier (\var{cart\_count}) et la
liste des categories pour la sidebar (\var{sidebar\_categories}), sans avoir
a repeter cette logique dans chaque vue.

% =============================================================================
\chapter{Moteur de recommandation}
% =============================================================================

\section{Vue d'ensemble de l'architecture}

Le moteur de recommandation est decompose en quatre modules Python, chacun
ayant un role distinct.

\begin{figure}[h]
\centering
\begin{tikzpicture}[
  box/.style={rectangle, draw=gold, thick, rounded corners=4pt,
              minimum width=3.5cm, minimum height=1cm,
              font=\small, align=center},
  arrow/.style={-Stealth, thick},
]
  \node[box, fill=darkbg!15] (hybrid)   at (4, 4)   {\textbf{hybrid.py}\\Orchestrateur};
  \node[box, fill=blue!10]   (collab)   at (0, 2)   {collaborative.py\\User-based + Item-based};
  \node[box, fill=green!10]  (content)  at (8, 2)   {content\_based.py\\TF-IDF + Cosinus};
  \node[box, fill=gray!15]   (cold)     at (4, 2)   {Cold Start\\Popularite};
  \node[box, fill=red!10]    (eval)     at (4, 0)   {evaluation.py\\P@K, R@K, F1, NDCG};

  \draw[arrow] (hybrid) -- (collab);
  \draw[arrow] (hybrid) -- (content);
  \draw[arrow] (hybrid) -- (cold);
  \draw[arrow] (hybrid) -- (eval);
\end{tikzpicture}
\caption{Architecture du moteur de recommandation}
\end{figure}

\section{Filtrage collaboratif (collaborative.py)}

\subsection{Construction de la matrice utilisateur-produit}

La fonction \fonction{build\_user\_item\_matrix} construit une matrice
$U \times P$ (utilisateurs $\times$ produits) a partir du queryset de notes.
Chaque cellule $M_{u,p}$ contient le score donne par l'utilisateur $u$ au
produit $p$, ou 0 si l'utilisateur n'a pas note ce produit.

\begin{lstlisting}[language=Python,caption={Construction de la matrice}]
matrix = np.zeros((len(user_ids), len(item_ids)))
for r in ratings_qs:
    matrix[user_idx[r.user_id]][item_idx[r.product_id]] = r.score
\end{lstlisting}

\subsection{Calcul de la similarite cosinus}

La similarite cosinus entre deux vecteurs $\mathbf{a}$ et $\mathbf{b}$ est :

\begin{equation}
\text{sim}(\mathbf{a}, \mathbf{b}) = \frac{\mathbf{a} \cdot \mathbf{b}}{||\mathbf{a}|| \cdot ||\mathbf{b}||}
\end{equation}

L'implementation utilise une normalisation matricielle avec NumPy, ce qui
permet de calculer toutes les similarites en une seule operation matricielle :

\begin{lstlisting}[language=Python,caption={Calcul matriciel de la similarite cosinus}]
def cosine_similarity_matrix(matrix):
    norms = np.linalg.norm(matrix, axis=1, keepdims=True)
    norms[norms == 0] = 1e-10  # evite la division par zero
    normalized = matrix / norms
    return normalized @ normalized.T
\end{lstlisting}

\subsection{Filtrage user-based}

L'algorithme identifie les \var{n\_similar} (par defaut 5) utilisateurs les
plus similaires a l'utilisateur cible. Pour chaque produit non encore note par
l'utilisateur cible, un score predit est calcule comme la moyenne ponderee par
similarite des notes des voisins :

\begin{equation}
\hat{r}_{u,p} = \frac{\sum_{v \in V} \text{sim}(u, v) \cdot r_{v,p}}{\sum_{v \in V} \text{sim}(u, v)}
\end{equation}

ou $V$ est l'ensemble des voisins ayant note le produit $p$.

\subsection{Filtrage item-based}

La meme logique est appliquee en transposant la matrice. La similarite est
maintenant calculee entre produits (colonnes de la matrice originale). Pour
chaque produit non note par l'utilisateur, le score predit est une moyenne
ponderee des notes de l'utilisateur sur les produits similaires :

\begin{equation}
\hat{r}_{u,p} = \frac{\sum_{q \in Q} \text{sim}(p, q) \cdot r_{u,q}}{\sum_{q \in Q} \text{sim}(p, q)}
\end{equation}

ou $Q$ est l'ensemble des produits notes par l'utilisateur similaires a $p$.

\section{Filtrage par contenu (content\_based.py)}

\subsection{Construction des descripteurs de produits}

Pour chaque produit, un descripteur textuel est construit en concatenant son
nom, sa categorie, sa description, ses tags et sa marque. Ce texte est mis en
minuscules et nettoye :

\begin{lstlisting}[language=Python,caption={Construction du descripteur}]
def build_product_features(products):
    features = []
    for p in products:
        text = f"{p.name} {p.category.name} {p.description} {p.tags} {p.brand}"
        features.append(text.lower().strip())
    return features
\end{lstlisting}

\subsection{Vectorisation TF-IDF}

Le vectoriseur \texttt{TfidfVectorizer} de scikit-learn transforme ces
descripteurs en vecteurs numeriques. Les stop words anglais sont supprimes,
et le vocabulaire est limite a 500 termes (\var{max\_features=500}).

TF-IDF (Term Frequency - Inverse Document Frequency) attribue a chaque terme
un poids refletant son importance relative dans le corpus :

\begin{equation}
\text{tfidf}(t, d) = \text{tf}(t, d) \times \log\frac{N}{\text{df}(t)}
\end{equation}

ou $\text{tf}(t, d)$ est la frequence du terme $t$ dans le document $d$, $N$
le nombre total de documents, et $\text{df}(t)$ le nombre de documents
contenant $t$.

\subsection{Deux modes d'utilisation}

Le module expose deux fonctions avec des usages distincts.

\fonction{content\_based\_recommendations} prend un produit cible et retourne
les produits les plus similaires selon la similarite cosinus entre vecteurs
TF-IDF. Elle est utilisee dans la page detail pour la section "produits
similaires".

\fonction{content\_based\_for\_user} calcule le profil de l'utilisateur comme
la moyenne des vecteurs TF-IDF de tous ses produits notes, puis retourne les
produits les plus proches de ce profil moyen. Elle est utilisee dans le systeme
hybride.

\begin{lstlisting}[language=Python,caption={Profil utilisateur par contenu}]
rated_indices = [product_ids.index(pid)
                 for pid in user_rated_product_ids if pid in product_ids]
user_profile = np.mean(tfidf_matrix[rated_indices].toarray(), axis=0)
similarities = cosine_similarity([user_profile], tfidf_matrix.toarray())[0]
\end{lstlisting}

\section{Systeme hybride (hybrid.py)}

\subsection{Principe de combinaison}

La fonction \fonction{hybrid\_recommendations} orchestre les trois approches
selon la formule :

\begin{equation}
\text{score\_final}(p) = \alpha \cdot \text{score\_collab}(p) + (1 - \alpha) \cdot \text{score\_contenu}(p)
\end{equation}

avec $\alpha = 0.5$ par defaut, donnant un poids egal aux deux approches.

\subsection{Construction du score collaboratif}

Les deux algorithmes collaboratifs sont executes chacun avec $2n$ resultats.
Un score de position (rang inverse) est attribue a chaque produit, puis les
scores des deux listes sont additionnes :

\begin{lstlisting}[language=Python,caption={Score collaboratif par rang}]
collab_scores = {}
for i, pid in enumerate(user_recs):
    collab_scores[pid] = collab_scores.get(pid, 0) + (len(user_recs) - i)
for i, pid in enumerate(item_recs):
    collab_scores[pid] = collab_scores.get(pid, 0) + (len(item_recs) - i)
\end{lstlisting}

\subsection{Normalisation}

Les scores collaboratifs et par contenu sont normalises independamment par
leur valeur maximale respective, ramenant les deux echelles sur $[0, 1]$ :

\begin{lstlisting}[language=Python,caption={Normalisation et fusion}]
max_c  = max(collab_scores.values(),  default=1) or 1
max_cb = max(content_scores.values(), default=1) or 1

for pid in all_ids:
    if pid in rated_ids:
        continue
    c_score  = collab_scores.get(pid, 0)  / max_c
    cb_score = content_scores.get(pid, 0) / max_cb
    final_scores[pid] = alpha * c_score + (1 - alpha) * cb_score
\end{lstlisting}

\subsection{Gestion du cold start}

Le cold start designe la situation ou un utilisateur n'a aucune note,
rendant les algorithmes collaboratifs inutilisables. Le systeme le detecte
en amont et bascule sur une heuristique de popularite :

\begin{lstlisting}[language=Python,caption={Gestion du cold start}]
if user_ratings.count() == 0:
    return cold_start_recommendations(all_products, n)
\end{lstlisting}

\fonction{cold\_start\_recommendations} retourne les produits tries par
\var{is\_featured} en premier, puis par note moyenne decroissante, puis par
nombre de notes. Ce mecanisme est egalement declenche si la combinaison hybride
ne produit aucun resultat.

\section{Evaluation (evaluation.py)}

Quatre metriques sont implementees. Toutes sont calculees en comparant la liste
de recommandations aux produits que l'utilisateur a notes $\geq 4/5$ (seuil de
pertinence).

\subsection{Precision@K}

Proportion de recommandations pertinentes parmi les $K$ premiers resultats :

\begin{equation}
\text{Precision@K} = \frac{|\text{recommandes}_{K} \cap \text{pertinents}|}{K}
\end{equation}

\subsection{Recall@K}

Proportion de produits pertinents effectivement retrouves dans les $K$ premiers
resultats :

\begin{equation}
\text{Recall@K} = \frac{|\text{recommandes}_{K} \cap \text{pertinents}|}{|\text{pertinents}|}
\end{equation}

\subsection{F1-Score}

Moyenne harmonique de la Precision et du Recall, penalisant les cas
desequilibres :

\begin{equation}
\text{F1} = \frac{2 \cdot \text{Precision@K} \cdot \text{Recall@K}}{\text{Precision@K} + \text{Recall@K}}
\end{equation}

\subsection{NDCG@K}

Le Normalized Discounted Cumulative Gain prend en compte la position des
elements pertinents dans la liste. Un produit pertinent en premiere position
contribue davantage qu'un produit pertinent en dixieme position :

\begin{equation}
\text{DCG@K} = \sum_{i=1}^{K} \frac{\mathbb{1}[p_i \in \text{pertinents}]}{\log_2(i + 1)}
\end{equation}

\begin{equation}
\text{NDCG@K} = \frac{\text{DCG@K}}{\text{IDCG@K}}
\end{equation}

ou $\text{IDCG@K}$ est le DCG ideal obtenu si tous les elements pertinents
etaient places en tete de liste.

\begin{lstlisting}[language=Python,caption={Implementation du NDCG@K}]
def ndcg_at_k(recommended_ids, relevant_ids, k=10):
    import math
    relevant_set = set(relevant_ids)
    top_k = recommended_ids[:k]
    dcg = 0.0
    for i, pid in enumerate(top_k):
        if pid in relevant_set:
            dcg += 1.0 / math.log2(i + 2)
    ideal_hits = min(len(relevant_set), k)
    idcg = sum(1.0 / math.log2(i + 2) for i in range(ideal_hits))
    return round(dcg / idcg, 3) if idcg > 0 else 0.0
\end{lstlisting}

Les metriques sont affichees dans l'interface uniquement quand l'utilisateur
a au moins 3 notes, garantissant que le calcul a du sens.

% =============================================================================
\chapter{Routes et vues}
% =============================================================================

\section{Table des routes}

\begin{table}[h]
\centering
\begin{tabularx}{\linewidth}{@{}llX@{}}
\toprule
\textbf{URL} & \textbf{Nom} & \textbf{Description} \\
\midrule
\texttt{/}                        & index          & Redirige vers landing ou home \\
\texttt{/landing/}                & landing        & Page d'accueil visiteurs \\
\texttt{/home/}                   & home           & Accueil utilisateur connecte \\
\texttt{/login/}                  & login          & Connexion \\
\texttt{/register/}               & register       & Inscription \\
\texttt{/email-sent/}             & email\_sent    & Confirmation envoi email \\
\texttt{/activate/<uid>/<token>/} & activate       & Activation de compte \\
\texttt{/logout/}                 & logout         & Deconnexion \\
\texttt{/products/}               & products       & Catalogue avec recherche et tri \\
\texttt{/products/<pk>/}          & product\_detail & Detail d'un produit \\
\texttt{/cart/}                   & cart           & Panier \\
\texttt{/cart/add/<pk>/}          & add\_to\_cart  & Ajout au panier (AJAX) \\
\texttt{/cart/update/<pk>/}       & update\_cart   & Modification quantite \\
\texttt{/cart/remove/<pk>/}       & remove\_from\_cart & Suppression (AJAX) \\
\texttt{/cart/checkout/}          & checkout       & Finalisation commande \\
\texttt{/profile/}                & profile        & Profil utilisateur \\
\texttt{/profile/update/}         & update\_profile & Mise a jour profil \\
\texttt{/recommendations/}        & recommendations & Page de recommandations \\
\texttt{/api/rate/<pk>/}          & api\_rate      & API JSON notation \\
\bottomrule
\end{tabularx}
\caption{Table complete des routes}
\end{table}

\section{Protection des acces}

Toutes les vues sont protegees par le decorateur \texttt{@login\_required},
a l'exception de \fonction{login\_view}, \fonction{register\_view},
\fonction{landing\_view}, \fonction{index\_view}, \fonction{email\_sent\_view}
et \fonction{activate\_view}. La redirection non authentifiee pointe vers
\texttt{/login/} comme defini dans \var{LOGIN\_URL}.

\section{Interactions AJAX}

Deux vues supportent les requetes AJAX via la detection de l'en-tete
\texttt{X-Requested-With: XMLHttpRequest} : \fonction{add\_to\_cart} retourne
le nouveau compteur du panier en JSON, et \fonction{remove\_from\_cart} retourne
le nouveau total et le nouveau compteur. L'endpoint \texttt{/api/rate/<pk>/}
est entierement JSON pour la notation depuis la page detail.

% =============================================================================
\chapter{Donnees de test}
% =============================================================================

\section{Commande seed\_data}

La commande \texttt{python manage.py seed\_data} peuple la base de donnees
avec un jeu de donnees representatif pour permettre l'utilisation immediate du
moteur de recommandation.

\subsection{Categories}

12 categories sont creees avec leurs icones, couleurs et slugs :
Electronique, Mode, Maison, Sport, Beaute, Livres, Cuisine, Jeux,
Voyage, Auto, Bijoux et Bebe.

\subsection{Produits}

Plus de 30 produits sont inseres, couvrant toutes les categories, avec :
\begin{itemize}
  \item Des descriptions detaillees integrant nom, caracteristiques et marque.
  \item Des tags en virgule-separes pour le filtrage TF-IDF.
  \item Des URLs Unsplash pour les images.
  \item Des prix et prix originaux permettant le calcul de reduction.
  \item Un stock aleatoire entre 10 et 200 unites.
\end{itemize}

\subsection{Utilisateurs de test}

9 utilisateurs sont crees (dont un compte demo) avec le mot de passe
\texttt{pass1234} ou \texttt{demo1234}. Un compte administrateur est gere
separement via \texttt{createsuperuser}.

\begin{table}[h]
\centering
\begin{tabular}{@{}lll@{}}
\toprule
\textbf{Email} & \textbf{Mot de passe} & \textbf{Role} \\
\midrule
demo@luxemart.com  & demo1234  & Compte demonstration \\
alice@luxemart.com & pass1234  & Utilisateur test \\
bob@luxemart.com   & pass1234  & Utilisateur test \\
clara@luxemart.com & pass1234  & Utilisateur test \\
\ldots             & pass1234  & (et 5 autres) \\
admin@luxemart.com & admin1234 & Superadministrateur \\
\bottomrule
\end{tabular}
\caption{Comptes de test disponibles}
\end{table}

\subsection{Notes generees}

Chaque utilisateur note 20 produits tires aleatoirement. La distribution des
scores est biaisee vers les bonnes notes pour simuler un comportement realiste :

\begin{lstlisting}[language=Python,caption={Distribution des notes dans le seed}]
score = random.choices([1, 2, 3, 4, 5], weights=[5, 10, 20, 35, 30])[0]
\end{lstlisting}

Cela genere environ 180 notes au total, suffisantes pour que les algorithmes
collaboratifs produisent des resultats significatifs.

% =============================================================================
\chapter{Interface utilisateur}
% =============================================================================

\section{Design et charte graphique}

L'interface adopte une esthetique sobre et luxueuse avec la palette suivante :

\begin{itemize}
  \item Or : \texttt{\#D4AF37} pour les elements d'accentuation.
  \item Fond sombre : \texttt{\#1a1a2e} pour les en-tetes et elements de mise
        en valeur.
  \item Blanc creme pour le contenu principal.
\end{itemize}

La typographie combine Cormorant Garamond pour les titres et Jost pour le
corps du texte, chargees depuis Google Fonts. L'interface supporte un mode
sombre et un mode clair, persiste dans le \texttt{localStorage} du navigateur.
La mise en page est responsive avec une sidebar retractable sur mobile.

\section{Templates principaux}

\begin{description}
  \item[\fichier{base.html}] Template parent contenant la navigation, la
    sidebar categories, le compteur de panier, et les blocs extensibles.
  \item[\fichier{home.html}] Affiche les categories, produits vedettes, mieux
    notes, et les recommandations personnalisees.
  \item[\fichier{products.html}] Catalogue avec barre de recherche, filtres
    par categorie et options de tri.
  \item[\fichier{product\_detail.html}] Page produit avec galerie, formulaire
    de note, formulaire de commentaire, et section produits similaires.
  \item[\fichier{cart.html}] Panier interactif avec mise a jour AJAX des
    quantites et du total.
  \item[\fichier{profile.html}] Historique des achats, notes donnees,
    recommandations personnelles et metriques d'evaluation.
  \item[\fichier{recommendations.html}] Page dediee aux recommandations avec
    affichage des quatre metriques d'evaluation.
  \item[\fichier{partials/product\_card.html}] Composant reutilisable affichant
    une carte produit (image, nom, prix, note, bouton panier).
\end{description}

% =============================================================================
\chapter{Limitations et perspectives}
% =============================================================================

\section{Limitations actuelles}

\subsection{Evaluation sur les memes donnees}

La principale limitation methodologique de l'evaluation est l'absence de
separation train/test. Les metriques sont calculees en comparant les
recommandations aux notes existantes de l'utilisateur, les memes notes qui ont
servi a construire le modele. Pour une evaluation rigoureuse, il faudrait cacher
une partie des notes avant de generer les recommandations, puis verifier si ces
notes cachees se retrouvent dans les resultats.

\subsection{Performance de la vectorisation}

La vectorisation TF-IDF est recalculee a chaque requete, sans mise en cache.
Avec un catalogue de grande taille, cette operation peut devenir un goulot
d'etranglement. Une amelioration serait de precalculer et stocker les matrices
de similarite.

\subsection{Stop words}

Le vectoriseur TF-IDF utilise des stop words anglais uniquement, alors que les
descripteurs de produits sont en francais. Certains termes tres frequents en
francais (articles, prepositions) pourraient diluer la qualite des vecteurs.

\section{Perspectives d'amelioration}

\begin{description}
  \item[Evaluation train/test] Implementer une separation des donnees en
    ensemble d'entrainement et ensemble de test pour des metriques plus fiables.
  \item[Stop words francais] Utiliser une liste de stop words francais dans
    le \texttt{TfidfVectorizer} ou passer par une librairie comme spaCy.
  \item[Mise en cache] Utiliser le cache Django (Redis) pour stocker les
    matrices de similarite precalculees et eviter leur recalcul a chaque requete.
  \item[Factorisation matricielle] Implementer SVD (Singular Value Decomposition)
    ou ALS (Alternating Least Squares) comme alternative aux methodes
    de voisinage pour ameliorer la qualite des predictions.
  \item[Dataset Amazon Reviews] Integrer le dataset Amazon Reviews 2023
    (mentionne dans le README) pour remplacer les donnees de seed synthetiques
    par des donnees reelles.
  \item[Ajustement du parametre alpha] Permettre a l'administrateur d'ajuster
    dynamiquement le parametre $\alpha$ du systeme hybride selon les metriques
    d'evaluation observees.
\end{description}

% =============================================================================
\chapter{Conclusion}
% =============================================================================

RecoShop est une application e-commerce complete integrant un moteur de
recommandation hybride fonctionnel. Le projet illustre la mise en oeuvre
pratique de plusieurs algorithmes classiques de systemes de recommandation :
filtrage collaboratif base utilisateur, filtrage collaboratif base article, et
filtrage par contenu TF-IDF.

L'architecture est proprement separee : chaque algorithme occupe son propre
module Python, le systeme hybride orchestre la combinaison dans un seul point
d'entree, et l'evaluation est completement independante du reste. Cette
modularite facilite l'ajout de nouveaux algorithmes ou le remplacement d'un
composant existant sans impacter le reste du code.

L'integration de quatre metriques d'evaluation (Precision@K, Recall@K, F1-Score
et NDCG@K) directement dans l'interface utilisateur offre une transparence
rare dans les applications e-commerce, permettant a chaque utilisateur de
comprendre la qualite des recommandations qui lui sont proposees.

Le deploiement est entierement automatise via un script de build et une
configuration par variables d'environnement, rendant l'application prete
pour un environnement de production sur Render avec PostgreSQL.

% =============================================================================
% Bibliographie
% =============================================================================
\begin{thebibliography}{9}

\bibitem{ricci2011}
Ricci, F., Rokach, L., \& Shapira, B. (2011).
\textit{Introduction to Recommender Systems Handbook}.
Springer.

\bibitem{sklearn}
Pedregosa, F. et al. (2011).
Scikit-learn: Machine Learning in Python.
\textit{Journal of Machine Learning Research}, 12, 2825--2830.

\bibitem{ndcg}
Jarvelin, K., \& Kekalainen, J. (2002).
Cumulated gain-based evaluation of IR techniques.
\textit{ACM Transactions on Information Systems}, 20(4), 422--446.

\bibitem{django}
Django Software Foundation. (2024).
\textit{Django Documentation, Version 5.2}.
\url{https://docs.djangoproject.com/}

\bibitem{tfidf}
Salton, G., \& Buckley, C. (1988).
Term-weighting approaches in automatic text retrieval.
\textit{Information Processing and Management}, 24(5), 513--523.

\end{thebibliography}

\end{document}